\documentclass[10pt]{article}
\usepackage[margin=20mm]{geometry}
\usepackage{amsmath,amssymb,mathtools,fontspec,unicode-math,xurl,hyperref}
\IfFontExistsTF{DejaVu Serif}{\setmainfont{DejaVu Serif}}{\setmainfont{Latin Modern Roman}}
\setmathfont{latinmodern-math.otf}
\setlength{\parindent}{0pt}\setlength{\parskip}{6pt}\sloppy
\title{Technical Verification of
Marked Length Spectrum Rigidity for Anosov Surfaces}
\author{AI-assisted technical review}
\date{}
\begin{document}
\maketitle
Abstract: The smooth rigidity argument is checked through its invariant-distribution extension, wavefront cone, period pairing, Torelli step, finite covers and conformal closing argument. No substantive mathematical defect is established in that inspected chain. The pairing normalization and the finite-cover separation have explicit supporting checks. A local tangent-orientation display needs a sign clarification without changing the operative formula. The integer regularity endpoint in the separate \(C^4\) extension remains a bounded Needs Evidence item; no counterexample to that extension is asserted. The reviewed file is an exact author version, not the journal version of record. Scientific adjudication remains pending.\par
\section*{Identity, version and inspected scope}
The paper is Colin Guillarmou, Thibault Lefeuvre and Gabriel P. Paternain, Marked length spectrum rigidity for Anosov surfaces, associated with Duke Mathematical Journal 174 (2025), 131–157, DOI 10.1215/00127094-2024-0022 (https://doi.org/10.1215/00127094-2024-0022). The exact reviewed document is arXiv:2303.12007v5 (https://arxiv.org/abs/2303.12007v5), dated 6 September 2024, with 22 PDF pages [paper]. Equivalence to the journal version has not been established. All locations below use this author PDF’s physical page numbers, which agree with its printed page numbers.\par
Actual source-image inspections cover all 22 pages across those roles. Reading and image inspection do not mean that every external theorem has been independently reproved. The principal inspected claims are Theorem 1.1, Lemmas 2.4–2.8, Propositions 3.1–3.7, Lemma 3.8 and the separate Remark 3.9/Lemma 3.10. Corollaries 1.2–1.3 and the unused higher-mode and algebra remarks have the narrower dependency coverage recorded below. No separate supplement or original computational repository was identified in bounded checks; their absence is not a mathematical defect. No numbered substantive Problem is manufactured.\par
\section*{Definitions and the smooth analytic chain}
The operative setting is a closed connected oriented surface with an Anosov geodesic flow. Its genus is at least two. The marked length spectrum is understood on nontrivial free homotopy classes; the trivial class can instead be assigned length zero if the codomain is adjusted. This harmless convention is needed for the positive-length notation on p. 1.\par
Write \(X\) for the geodesic field, \(V\) for positive fiber rotation and \(H=[V,X]\). The dual coframe \((\lambda,\beta,\psi)\) satisfies \[d\lambda=\psi\wedge\beta,\qquad
 \mu=-\lambda\wedge d\lambda=\lambda\wedge\beta\wedge\psi.\] The operators \(\eta_\pm=(X\mp iH)/2\) change Fourier degree by \(\pm1\). In isothermal coordinates \(g=e^{2\rho}|dz|^2\), \[\eta_-\bigl(ae^{in\theta}\bigr)
 =e^{-(n+1)\rho}\bar\partial(ae^{n\rho})e^{i(n-1)\theta}.\] This verifies the holomorphic-bundle identification including its metric factor. For a distribution \(f\in\mathcal A_+\), the \(-1\) and zero modes of \(Xf=0\) give \(\eta_-f_0=\eta_-f_1=0\). Base elliptic regularity makes these coefficients smooth; compactness makes \(f_0\) constant. No initial smoothness of an arbitrary distribution’s Fourier coefficients is presumed.\par
The first-mode extension has the required hypotheses. In the exact consulted version of Paternain–Salo–Uhlmann [psu], Theorem 5.6, p. 18, supplies \[w=\sum_{k\ge1}w_k\in L_x^2H_\theta^{-1},\qquad
 Xw=0,\qquad w_1=a_1\] for any \(\eta_-a_1=0\) on a smooth Anosov surface. It imposes no negative-curvature or nonhyperelliptic restriction. With the paper’s adjoint convention \(\pi_{1*}\pi_1^*=2\pi\operatorname{Id}\), set \(a_1=(2\pi)^{-1}\pi_1^*\omega\). This produces \(\pi_{1*}w=\omega\). A choice on a finite complex basis gives a complex-linear right inverse \(e_1\), with no uniqueness requirement. The stronger higher-mode remark can use the solenoidal extension of [gui], Corollary 3.8, together with its surface injectivity Theorem 3.12 and a parity/truncation argument. Its kernel-orthogonality hypothesis is not silently dropped, and that stronger result is not needed for the main first-mode proof.\par
For the Anosov splitting write \(Y^s=H+r_-V\), \(Y^u=H+r_+V\), with \(r_+>r_-\). The inequality follows from the positive transverse Jacobi coefficients and their Wronskian; it does not assume curvature is everywhere negative. In the characteristic plane \(\xi(X)=0\), put \[\xi_s=\psi-r_-\beta,\qquad \xi_u=\psi-r_+\beta.\] Both evaluate to one on \(V\). Their lifted-flow factors are positive and have opposite exponential behaviour. A negative coefficient on either ray eventually exits \(\{\xi(V)\ge0\}\) in one time direction. Consequently its maximal two-sided invariant subset is precisely the positive cone generated by these rays. This supplies the central step in Lemma 2.5. Every nonzero covector in that cone has positive vertical component, so fiber push-forward makes every Fourier coefficient smooth.\par
An exact smooth time-preserving conjugacy preserves \(X\), the individual stable and unstable bundles, and therefore \(\lambda\) and contact volume. The leaf-space circle on the universal cover has degree \(+1\) under the identity-isotopy condition in Proposition 3.2. It fixes the transverse orientation; contact orientation then fixes the two individual ray signs. Thus pullback carries the positive cone to the positive cone, not its negative. For \(f'=\phi^*f\), the negative Szegö projection is smooth and \[X_2(S^-f')=\eta_+f'_0+\eta_+f'_{-1}.\] Only degrees zero and one occur. Degree-one tensor tomography makes \(S^-f'=f'_0\), eliminating all negative modes. The inverse conjugacy gives the inverse map on \(\mathcal A_+\). The needed Szegö statement includes its horizontal, zero-vertical kernel directions; these do not intersect the strict positive cone. The argument does not require a falsely diagonal-only kernel description.\par
\section*{Pairing, marking and the conformal conclusion}
\subsection*{The constant and the local sign clarification}
Equation reference: eq:pair. At angle \(\theta\) over an oriented geodesic, the coordinate tangent is \(T_\tau=\cos\theta\,X-\sin\theta\,H\). Thus the consistent orientation is \((\cos\theta\,\lambda-\sin\theta\,\beta)\wedge\psi\). The two corresponding plus signs preceding Lemma 2.7 on p. 11 are local display slips: that printed form evaluates to \(\cos(2\theta)\) on the actual coordinate basis. With the corrected signs, \[d\lambda=\sin\theta\,d\tau\wedge d\theta,\qquad
 \int_0^{2\pi}e^{ik\theta}\sin\theta\,d\theta
 =i\pi(\delta_{k1}-\delta_{k,-1}).\] For \(f\in\mathcal A_+\) the negative first mode vanishes, and the adjoint normalization therefore gives exactly the formula actually used: \[
 2i\int_{S^1\gamma}f\,d\lambda
 =-\int_{[\gamma]}\pi_{1*}f .\] The final sign and factor in Lemma 2.8 survive the clarification. Ten exact symbolic controls, including modes \(-4,\ldots,4\), support this calculation; the displayed Fourier identity is the all-integer analytic argument.\par
The conormal to the parametrized circle bundle annihilates \(V\), whereas the admissible wavefront cone has strictly positive vertical component. Thus its distributional pullback exists. Self-intersecting geodesics are handled as immersed cycles with multiplicity, or one can use simple representatives of a symplectic homology basis. Passing from smooth functions to distributions uses wavefront-controlled regularization, not unrestricted weak convergence.\par
\subsection*{Why homology suffices}
The form \(f\,d\lambda\) is closed because \(Xf=0\). For a direct smooth-category justification of the de Rham step on pp. 16–17, choose an auxiliary smooth metric, its harmonic projector \(\mathcal{H}\) and Green operator \(G\). For a closed current \(T\), \[T=\mathcal{H}T+d(d^*GT),\qquad
 \operatorname{WF}(d^*GT)\subseteq\operatorname{WF}(T).\] The identity follows from Hodge decomposition and \(dT=0\); the wavefront inclusion follows from the pseudodifferential nature of \(d^*G\). It gives a smooth closed representative plus a primitive whose traces are defined on both relevant cycles. Stokes’ formula then makes the pairing depend only on their homology. Together with \(\phi^*\lambda_1=\lambda_2\), this proves that \[F=\pi_{1*}\phi^*e_1:
 H^0_{J_1}(\Sigma,K)\longrightarrow H^0_{J_2}(\Sigma,K)\] preserves all periods. A holomorphic one-form with zero periods is zero, so \(F\) is injective and hence an isomorphism in genus \(g\). Normalized \(a\)-period bases therefore give identical period matrices.\par
Torelli yields an unmarked biholomorphism at this point. The paper correctly retains the mapping-class ambiguity. The exact author version of Alvarez–Brum–Martínez–Potrie with Wolff’s appendix [abmp], Theorem C on pp. 37–38, applies to every simple closed geodesic, including a separating one. It gives a finite cover with a degree-one lift of that curve and all other simple closed geodesics longer than any prescribed \(K\).\par
If two marked hyperbolic metrics differ, choose a simple class \(\alpha\) with unequal lengths. Compactness gives a uniform ratio bound \(C\) for the lengths of all free homotopy classes. Choose \(K/C>\max(\ell_1(\alpha),\ell_2(\alpha))\). The isolated lift is then the systole for both lifted metrics, with different lengths. The lifts cannot lie in the same mapping-class orbit. The comparison must include nonsimple projected curves; the proof explicitly does. Applying the period/Torelli conclusion on every finite cover proves Proposition 3.7.\par
Finally suppose the marked conformal representatives satisfy \(g_2=e^{2f}g_1\). Contact conjugacy gives equality of area, \(A_1=A_2\). Each \(g_1\)-closed geodesic has \(g_2\)-length at least the \(g_2\)-minimum in its free homotopy class, which equals its \(g_1\)-length. Periodic-orbit approximation of Liouville measure yields \[\int_\Sigma e^f\,dA_1\ge A_1,\qquad
 \left(\int_\Sigma e^f\,dA_1\right)^2
 \le A_1\int_\Sigma e^{2f}\,dA_1=A_1A_2=A_1^2.\] Equality in Cauchy–Schwarz makes \(e^f\) constant, and equal area makes that constant one. This checks the final conformal closing step under the established conjugacy and standard Anosov periodic-approximation interface.\par
\section*{Finite regularity and remaining dependency limits}
Lemma 3.10, pp. 19–20, has a valid endpoint trace estimate under its stated conormal wavefront condition. In local coordinates split Fourier space at \(|\xi_1|=A|\xi'|\). Inside the cone, the integration interval has length \(O(\langle\xi'\rangle)\) and \(\langle\xi\rangle\asymp\langle\xi'\rangle\). Cauchy–Schwarz gives the \(H^s\to H^{s-1/2}\) estimate for every real \(s\); the rapidly decreasing normal-cone remainder has smooth trace. The weight comparison is made while retaining that cone restriction. This avoids an invalid unrestricted weighted-integral comparison.\par
A distinct endpoint remains to be completed in Remark 3.9. The exact consulted Gogolev–Rodriguez Hertz version [gr] states conjugacy regularity \(C^{r_*}\), with \(r_*=r\) for nonintegral \(r\) and \(C^{r_*}=C^{r-1,\mathrm{Lip}}\) at an integer. For a \(C^4\) metric its flow is \(C^3\), so this citation directly gives \(C^{2,\mathrm{Lip}}\), whereas the remark says \(C^3\). This difference alone is not a counterexample to either regularity improvement or the \(C^4\) rigidity conclusion.\par
The weaker endpoint may suffice: it has the relevant bounded weak derivatives, and \(H^{-1}\) distributional pullback is well-defined. A finite Sobolev wavefront argument can split off an \(H^2\) remainder, whose trace is \(H^{3/2}\); the main admissible trace is \(H^{-3/2}\) and can be paired with \(C^2\) coefficients. However, this review does not complete the required endpoint microlocal covariance and propagation under \(C^{2,\mathrm{Lip}}\) coordinate changes, nor the entire low-regularity Pestov argument. The requested completion is a precise endpoint lemma and its application, or an independent stronger conjugacy result. This is retained as Needs Evidence, with the smooth result separately supported. No unsupported Minor or Major finding is assigned.\par
Corollary 1.2 uses positive Livšic and the geodesic-current volume comparison argument cited by the authors. The ordering and equality implications were traced at their interfaces; the full Croke–Dairbekov/Lopes–Thieullen proofs were not independently reproduced. A bounded author-site retrieval returned HTTP 403. Corollary 1.3 concerns a smooth, time-preserving conjugacy. The center of the unit-tangent-bundle fundamental group identifies the fiber subgroup, so the conjugacy induces a base-surface automorphism. After realizing its mapping class and applying Theorem 1.1, one reduces to a self-conjugacy with trivial base marking. The universal-cover endpoint description fixes each oriented geodesic; on a dense orbit, time preservation fixes a single shift, and continuity gives the same shift everywhere. Time reversal is outside this corollary’s hypothesis.\par
Classical Torelli, the orbit-equivalence construction, the complete smooth-conjugacy and tensor-tomography foundations, and the unused general-geometry algebra statement remain external inputs. Specific extension and finite-cover interfaces were consulted in exact author versions, without asserting equality to their later journal files.\par
\section*{Assessment and state}
No substantive defect is established in the inspected smooth main argument. The local orientation signs are clarified with the operative pairing unchanged. The \(C^4\) endpoint and the expressly bounded external proof coverage remain visible limitations.\par
Dimension | Recorded state\par
Issue assessment | No Issue Found within the inspected smooth scope.\par
Numbered Problems | None established.\par
Separate unresolved item | Needs Evidence: the \(C^4\) integer-regularity endpoint and its microlocal use.\par
Repair verification | Local sign clarification checked; no full low-regularity repair certified.\par
Author adoption | Not established; no author response supplied.\par
Scientific adjudication | Pending; human-review list empty.\par
Independent acceptance | Not performed.\par
9 C. Guillarmou, T. Lefeuvre and G. P. Paternain, Marked length spectrum rigidity for Anosov surfaces, arXiv:2303.12007v5 (https://arxiv.org/abs/2303.12007v5), 6 September 2024. Exact 22-page author version reviewed. Journal identity: Duke Math. J. 174 (2025), 131–157; doi:10.1215/00127094-2024-0022. G. P. Paternain, M. Salo and G. Uhlmann, Spectral rigidity and invariant distributions on Anosov surfaces, arXiv:1208.4943v3 (https://arxiv.org/abs/1208.4943v3), 28 April 2014. Theorems 1.5 and 5.5–5.6 and the stated Fourier interfaces. C. Guillarmou, Invariant distributions and X-ray transform for Anosov flows, arXiv:1408.4732v2 (https://arxiv.org/abs/1408.4732v2), 21 December 2015. Corollary 3.8 and Lemma 3.10/Theorem 3.12 interfaces. S. Alvarez, J. Brum, M. Martínez and R. Potrie, Topology of leaves for minimal laminations by hyperbolic surfaces, with an appendix by the authors and M. Wolff, arXiv:1906.10029v2 (https://arxiv.org/abs/1906.10029v2), 1 February 2022. Appendix A, Theorem C, pp. 37–38. A. Gogolev and F. Rodriguez Hertz, Smooth rigidity for 3-dimensional volume preserving Anosov flows and weighted marked length spectrum rigidity, arXiv:2210.02295v2 (https://arxiv.org/abs/2210.02295v2). Theorem 1.1 and the integer regularity convention.\par
\section*{Supporting checks}
The reviewed canonical PDF SHA256 is \[\texttt{021434e951063ed64d76d620d1bc3523}\] \[\texttt{93cbaffc386d3002718c97bf053967d9}.\] The matching verification directory preserves the original manifest, text/page correspondence, analytical notes, exact symbolic controls, all actual source-page observations, external retrieval limits, compiler commands and dependencies, and final report-page inspection. Author source files were available but were neither executed nor compiled. The current report was compiled from its own exact TeX.\par
The root role read and integrated the complete paper and wrote this report. No independent UI observation or backend equivalence is asserted. Scientific decision is pending and `human\_reviews[]` remains empty.\par
\end{document}
